\documentclass[10pt,a4paper]{amsart}
\usepackage[margin=19mm]{geometry}
\usepackage{amsmath,amssymb,mathtools}
\usepackage[scr=rsfs,cal=euler]{mathalfa}
\usepackage{xfrac}
\usepackage[unicode,hidelinks,bookmarks=false]{hyperref}
\newcommand{\ie}{i.e.\ }
\newcommand{\cf}{cf.\ }
\newcommand{\resp}{respectively\ }
\newcommand{\Subsection}{\subsection}
\providecommand{\nobreakdash}{\nobreak\hbox{-}}
\setlength{\emergencystretch}{2em}
\multlinegap0pt
\usepackage{bm}
\usepackage{bbm}
\usepackage{dsfont}
\usepackage{xspace}
\usepackage{cancel}
\usepackage{mathtools}
\usepackage{cleveref}
\usepackage{amsthm}
\makeatletter
\@ifundefined{counter}{\newtheorem{counter}{Counter}}{}
\@ifundefined{lemma}{\newtheorem{lemma}[counter]{Lemma}}{}
\@ifundefined{proposition}{\newtheorem{proposition}[counter]{Proposition}}{}
\makeatother
\newcommand{\E}{\mathbb{E}}
\newcommand{\PP}{\mathbb{P}}
\newcommand{\R}{\mathbb{R}}
\newcommand{\U}{\mathbb{U}}
\newcommand{\ball}[3]{B_{#3}(#1, #2)} %
\newcommand{\by}{\mathbf{y}}
\newcommand{\cB}{\mathcal{B}}
\newcommand{\cC}{\mathcal{C}}
\newcommand{\cD}{\mathcal{D}}
\newcommand{\cE}{\mathcal{E}}
\newcommand{\cF}{\mathcal{F}}
\newcommand{\cN}{\mathcal{N}}
\newcommand{\ls}{\lesssim}
\newcommand{\one}{\mathbb{1}}
\newcommand{\talpha}{\tilde{\alpha}}
\newcommand{\teta}{\tilde{\eta}}
\title[The Geometry of Cochains on Sampled Vietoris-Rips Complexes]{The Geometry of Cochains on Sampled Vietoris-Rips Complexes}
\author{Darrick Lee, Kelly Maggs}
\date{Source: arXiv:2608.11137v2 (CC BY 4.0)}
\begin{document}
\maketitle

\noindent\emph{Source and licence.} The Geometry of Cochains on Sampled Vietoris-Rips Complexes. Version arXiv:2608.11137v2 (CC BY 4.0), distributed under the Creative Commons Attribution 4.0 International licence, as recorded in the arXiv licence metadata for this exact version and shown on the abstract page https://arxiv.org/abs/2608.11137v2. Full rights notice: licenses/arxiv-2608.11137-CC-BY-4.0.txt.\par\smallskip

\noindent\textit{Editorial scope.} Counted unit: the Proposition whose source label is \texttt{prop:off\_diagonal\_covering\_variance} in the article (source0.tex lines 2175--2183, source label \texttt{prop:off\_diagonal\_covering\_variance}), delivered together with the article's own complete original proof of it (source0.tex lines 2184--2221, one \texttt{proof} environment of 38 lines). No mathematics is written, repaired or re-derived: statement and proof are the article's own text line for line. Required context retained uncounted: lem:concentration\_D\_k\_support (lines 1561--1570); lem:codiff\_pointwise\_variance (lines 2129--2136); lem:codiff\_covering\_bound (lines 2165--2170). Every definition, display and label of those lines is retained unchanged. Omitted from this package and not redistributed: 1--1560, 1571--2128, 2137--2164, 2171--2174, 2222--2419. Nothing inside a retained excerpt is omitted or altered: every retained body is delivered exactly as the article's lines are written, so no omission receipt of any kind accompanies this package. Citations: the retained excerpts carry no citation command at all, so no bibliography entry of the article is transcribed and no reference is redistributed. Reference adaptation: none was needed and none was made; every reference inside the delivered unit resolves inside it, so no unresolved reference or citation is produced. Printed slips kept literal: no printed word, symbol, hypothesis or number of the counted statement or of its proof was altered. Layout and class: the article's own class is \texttt{amsproc} with packages including mathpazo, fontenc, geometry, comment, caption, scalerel, nicefrac, xy, tikz-cd, tikz, rotating, setspace, bbm, listings; the wrapper is the lane's pinned \texttt{amsart} editorial wrapper at 10pt on A4, which restates the theorem-like environments the retained text needs and adds the article's own macro definitions that the retained text needs (\texttt{\textbackslash E}, \texttt{\textbackslash PP}, \texttt{\textbackslash R}, \texttt{\textbackslash U}, \texttt{\textbackslash ball}, \texttt{\textbackslash by}, \texttt{\textbackslash cB}, \texttt{\textbackslash cC}, \texttt{\textbackslash cD}, \texttt{\textbackslash cE}, \texttt{\textbackslash cF}, \texttt{\textbackslash cN}, \texttt{\textbackslash ls}, \texttt{\textbackslash one}, \texttt{\textbackslash talpha}, \texttt{\textbackslash teta}), with extra wrapper packages (bm, bbm, dsfont, xspace, cancel, mathtools, cleveref). The delivered numbering is the unit's own. Assets: no retained span contains a figure environment, a graphics-inclusion command, a PDF-page inclusion, a table or a TikZ picture; no image file is copied into this package, the package declares no asset dependency and no image file is redistributed with it. Licence: version arXiv:2608.11137v2 of this article is distributed under the Creative Commons Attribution 4.0 International licence, as recorded in the arXiv licence metadata for this exact version and shown on the abstract page https://arxiv.org/abs/2608.11137v2. A full rights notice is delivered at licenses/arxiv-2608.11137-CC-BY-4.0.txt. Characters: every retained excerpt is pure ASCII, so the delivered engine, XeLaTeX with its default Latin Modern text font, reports no missing character.
    \begin{lemma} \label{lem:concentration_D_k_support}
        Let $\one_{D^k_\epsilon(M)} : M^{k+1} \to \R$ be the indicator function on $D^k_\epsilon(M)$, which is symmetric. For a sufficiently large constant $C$, define
        \begin{align}
            \cE_\epsilon \coloneqq \left\{ \U_n[\one_{D^k_\epsilon(M)}] \leq C\epsilon^{km}\right\},
        \end{align}
        Then,
        \begin{align}
            \PP[\cE_\epsilon^c] \leq 4 \exp \left(-cn\epsilon^{km}\right).
        \end{align}
    \end{lemma}

\begin{lemma} \label{lem:codiff_pointwise_variance}
    For $\alpha\in\pi^*\cB_M^k$, we have
    \begin{align}
        \PP[|F_n^\alpha-\E[F_n^\alpha]|>t]
        \leq 4\exp\left(\frac{-cnt^2}{A\epsilon^{-2}+B\epsilon^{-m(k+1)-2}t}\right),
    \end{align}
    where $A,B>0$ depend only on $M,N,k$, and $\theta$.
\end{lemma}

\begin{lemma} \label{lem:codiff_covering_bound}
    Let $\cN_\gamma$ be the smallest number of $\cC^0(M)$-balls of radius $\gamma>0$, with centers in $\cB_M^k$, required to cover $\cB_M^k$. Then
    \begin{align}
        \log \cN_\gamma\ls \gamma^{-m/2}.
    \end{align}
\end{lemma}

\begin{proposition} \label{prop:off_diagonal_covering_variance}
    There is a constant $C_0>0$, depending only on the fixed data, such that the following holds. Let $p\in(0,1)$ and suppose $C_0\log(8/p)\leq n\epsilon^{(k+1)m}$. With probability at least $1-p$, we have
    \begin{align}
        \sup_{\alpha\in\pi^*\cB_M^k}|F_n^\alpha-\E[F_n^\alpha]|
        \ls \sqrt{\frac{\gamma^{-m/2}+\log(1/p)}{n\epsilon^2}}
        +\frac{\gamma^{-m/2}+\log(1/p)}{n\epsilon^{m(k+1)+2}}
        +\gamma\epsilon^{-2}.
    \end{align}
\end{proposition}

\begin{proof}
    Let $\{\ball{\cC^0(M)}{\teta^r}{\gamma}\}_{r=1}^{\cN}$ be a collection of $\cC^0(M)$-balls which cover $\cB_M^k$, with centers $\teta^r\in\cB_M^k$. Let $\eta^r=\pi^*\teta^r$. By a union bound and~\Cref{lem:codiff_pointwise_variance},
    \begin{align} \label{eq:codiff_union_variance}
        \PP\left[\sup_{r\in[\cN]}|F_n^{\eta^r}-\E[F_n^{\eta^r}]|>t\right]
        \leq 4\cN\exp\left(\frac{-cnt^2}{A\epsilon^{-2}+B\epsilon^{-m(k+1)-2}t}\right).
    \end{align}
    Let $\alpha=\pi^*\talpha\in\pi^*\cB_M^k$. Choose $\eta^r$ such that $\|\talpha-\teta^r\|_{\cC^0(M)}<\gamma$. Then, by the definition of $A_\epsilon$ and the lower bound $\kappa_{\theta,\epsilon}^{k-1}(\by)\geq c_\theta$ on the support of $\kappa_{\theta,\epsilon}^{k}(x,\by)$,
    \begin{align}
        |A^\alpha_\epsilon(x,\by)-A^{\eta^r}_\epsilon(x,\by)|\ls \gamma\epsilon^{k-m-2}\one_{D^k_\epsilon(M)}(x,\by).
    \end{align}
    If $\cD_{F,\epsilon}(M)$ denotes the support of $h^F_\epsilon$, then $\cD_{F,\epsilon}(M)\subset D^{k+1}_{2\epsilon}(M)$. Hence
    \begin{align}
        |h^F_\epsilon(\alpha)-h^F_\epsilon(\eta^r)|\ls \gamma\epsilon^{-m(k+1)-2}\one_{D^{k+1}_{2\epsilon}(M)}.
    \end{align}
    Taking expectations gives
    \begin{align}
        \sup_{\alpha\in\pi^*\cB_M^k}|\E[h^F_\epsilon(\alpha)-h^F_\epsilon(\eta^r)]|\ls \gamma\epsilon^{-2}.
    \end{align}
    By the same argument as~\Cref{lem:concentration_D_k_support}, applied to $D^{k+1}_{2\epsilon}(M)$, there is an event
    \begin{align}
        \cF_\epsilon \coloneqq \left\{\U_n[\one_{D^{k+1}_{2\epsilon}(M)}]\leq C\epsilon^{(k+1)m}\right\}
    \end{align}
    such that $\PP[\cF_\epsilon^c]\leq 4\exp(-cn\epsilon^{(k+1)m})$. This is distinct from $\cE_\epsilon$, which controls the $D^k_\epsilon(M)$ support. On $\cF_\epsilon$,
    \begin{align}
        \sup_{\alpha\in\pi^*\cB_M^k}|\U_n[h^F_\epsilon(\alpha)-h^F_\epsilon(\eta^r)]|\ls \gamma\epsilon^{-2}.
    \end{align}
    Combining this with~\eqref{eq:codiff_union_variance} and the triangle inequality gives
    \begin{align}
        \PP\left[\sup_{\alpha\in\pi^*\cB_M^k}|F_n^\alpha-\E[F_n^\alpha]|>t+C\gamma\epsilon^{-2}\right]
        \leq 4\cN\exp\left(\frac{-cnt^2}{A\epsilon^{-2}+B\epsilon^{-m(k+1)-2}t}\right)+4\exp(-cn\epsilon^{(k+1)m}).
    \end{align}
    Using~\Cref{lem:codiff_covering_bound} and the standard Bernstein inversion with
    \begin{align}
        t\gtrsim \sqrt{\frac{\gamma^{-m/2}+\log(1/p)}{n\epsilon^2}}
        +\frac{\gamma^{-m/2}+\log(1/p)}{n\epsilon^{m(k+1)+2}}
    \end{align}
    proves the result.
\end{proof}

\end{document}
